% --- Archon physics formalization source begin ---
% archon:chemistry
% archon:covers IChO2026Problems/problem_icho_2026_t3_a7.lean
% archon:problem-id icho_2026_t3
% archon:part-id T3-A7

\chapter{Icho Chemistry problem icho\_2026\_t3\_a7}
\label{ch:IChO2026Problems_problem_icho_2026_t3_a7}

\paragraph{Problem source.}
\#\# Shared official problem context
T3. Into Reticular Chemistry 7\% Covalent organic frameworks (COFs) are a class of polymers that form two- or three-dimensional structures through reactions between organic precursors, resulting in covalent bonds to afford porous, crystalline materials. Two examples of honeycomb type 2D-COFs formed by boronate ester and boroxine rings (repeat units indicated using dashed lines) are shown in the figure below. The black and red bold lines represent different building blocks in all COF topology figures. 3.1 Determine the empirical formula of COF-1 and calculate the mass percentage (\%, to two decimal places) of carbon in COF-1. 3.2 Calculate the internal diameter, d, in Å, of the honeycomb for COF-2, with the following assumptions: • C–C/C=C (arenes) = 1.39 Å • C–B = 1.56 Å • B–O = 1.38 Å • Note, the diameter of a circle, 𝑑, inscribed inside a hexagon, is given by 𝑑= √3𝑎, where 𝑎is the side length of the hexagon. • Neglect the width of the linkers. The structures shown below can form 2D and 3D COFs. For this question consider only formation of COFs with boronate esters, imines, and C=C bonds through condensation reactions. A–D represent compound classes with different functional groups. Tetragonal means Tetragonal, Hexagonal means Hexagonal, Trigonal means Trigonal, Kagome means Kagome, Tetrahedral means Tetrahedral. 3.3 Fill all the table cells with either new examples of monomer combinations that give the topologies shown above or with “XXX” if you are certain that no additional combinations satisfy that topology. Both the combinations and the “XXX” entries will be graded. Incorrect answers (both combinations and “XXX” entries) will be penalised. Blank table cells will not be graded. Tetragonal 1 Tetragonal 2 Hexagonal 1 Hexagonal 2 Trigonal Kagome Tetrahedral A2 + B3 E1 + D2 C2 + D4 Whilst imine-based COFs are easily synthesised, they are also vulnerable to hydrolysis under acidic and basic conditions. The addition of hydroxy groups to the aldehyde or amine monomers can overcome this. IR spectrum of COF-3 contained stretches corresponding to the imine functional group, but COF-4 did not. COF-4 contains a different type of C=N bonds, and elemental analysis revealed a loss of six additional hydrogen atoms per repeat unit relative to COF-3 as the only change. COF-5 irreversibly isomerises to COF-6, which did not contain imine stretches in its IR spectrum. 3.4 Draw the repeat units for COFs 3–6. Indicate the edges of each repeat unit with a dashed line as shown for COF-1 below. Macrocycle X was synthesised by post-synthetic modification of COF-7 (C2 + D4) followed by degradation. 3.5 Draw the smallest repeat unit of macrocycle X. π-π interactions result in different stacking modes of COF layers, some of which are shown below. Depending on stacking geometry, the interaction energy can be different as shown below. 𝐸/ kJ mol−1 Interaction geometry benzene (X = CH) – benzene (b-b) −7.9 −12.6 benzene-triazine (X = N) (b-t) −49.8 −55.6 triazine–triazine (t-t) −6.7 −16.7 COF-8, synthesised from (E1 + D2), is notable for its remarkable stability. The AA' slightly shifted mode of a COF-8 bilayer has a π-π stacking interaction energy of –67.1 kJ mol−1 between two layers of one repeat unit. 3.6 Calculate the π-π stacking energy between two layers of one repeat unit of COF-8, for the AA, AB and AB' arrangements. Assume π-π stacking happens only between aromatic units. COF-9, obtained from COF-8, adsorbs UO2+ 2 ions and loses its fluorescence. In an experiment, a suspension of 5.000 mg COF-9 in 19.90 mg dm−3 UO2+ 2 solution (200.0 mL) was filtered after equilibrium was reached. The final concentration of UO2+ 2 was 9.225 mg dm−3.

\#\# Current subquestion T3-A7
Calculate the equilibrium absorption capacity at a given temperature (𝑞𝑒) in mg g−1 of COF-9. Indicate the number of UO2+ 2 ions absorbed per pore. Assume that all uranium exists as UO2+ 2 ions.

\paragraph{Shared problem context.}
T3. Into Reticular Chemistry 7\% Covalent organic frameworks (COFs) are a class of polymers that form two- or three-dimensional structures through reactions between organic precursors, resulting in covalent bonds to afford porous, crystalline materials. Two examples of honeycomb type 2D-COFs formed by boronate ester and boroxine rings (repeat units indicated using dashed lines) are shown in the figure below. The black and red bold lines represent different building blocks in all COF topology figures. 3.1 Determine the empirical formula of COF-1 and calculate the mass percentage (\%, to two decimal places) of carbon in COF-1. 3.2 Calculate the internal diameter, d, in Å, of the honeycomb for COF-2, with the following assumptions: • C–C/C=C (arenes) = 1.39 Å • C–B = 1.56 Å • B–O = 1.38 Å • Note, the diameter of a circle, 𝑑, inscribed inside a hexagon, is given by 𝑑= √3𝑎, where 𝑎is the side length of the hexagon. • Neglect the width of the linkers. The structures shown below can form 2D and 3D COFs. For this question consider only formation of COFs with boronate esters, imines, and C=C bonds through condensation reactions. A–D represent compound classes with different functional groups. Tetragonal means Tetragonal, Hexagonal means Hexagonal, Trigonal means Trigonal, Kagome means Kagome, Tetrahedral means Tetrahedral. 3.3 Fill all the table cells with either new examples of monomer combinations that give the topologies shown above or with “XXX” if you are certain that no additional combinations satisfy that topology. Both the combinations and the “XXX” entries will be graded. Incorrect answers (both combinations and “XXX” entries) will be penalised. Blank table cells will not be graded. Tetragonal 1 Tetragonal 2 Hexagonal 1 Hexagonal 2 Trigonal Kagome Tetrahedral A2 + B3 E1 + D2 C2 + D4 Whilst imine-based COFs are easily synthesised, they are also vulnerable to hydrolysis under acidic and basic conditions. The addition of hydroxy groups to the aldehyde or amine monomers can overcome this. IR spectrum of COF-3 contained stretches corresponding to the imine functional group, but COF-4 did not. COF-4 contains a different type of C=N bonds, and elemental analysis revealed a loss of six additional hydrogen atoms per repeat unit relative to COF-3 as the only change. COF-5 irreversibly isomerises to COF-6, which did not contain imine stretches in its IR spectrum. 3.4 Draw the repeat units for COFs 3–6. Indicate the edges of each repeat unit with a dashed line as shown for COF-1 below. Macrocycle X was synthesised by post-synthetic modification of COF-7 (C2 + D4) followed by degradation. 3.5 Draw the smallest repeat unit of macrocycle X. π-π interactions result in different stacking modes of COF layers, some of which are shown below. Depending on stacking geometry, the interaction energy can be different as shown below. 𝐸/ kJ mol−1 Interaction geometry benzene (X = CH) – benzene (b-b) −7.9 −12.6 benzene-triazine (X = N) (b-t) −49.8 −55.6 triazine–triazine (t-t) −6.7 −16.7 COF-8, synthesised from (E1 + D2), is notable for its remarkable stability. The AA' slightly shifted mode of a COF-8 bilayer has a π-π stacking interaction energy of –67.1 kJ mol−1 between two layers of one repeat unit. 3.6 Calculate the π-π stacking energy between two layers of one repeat unit of COF-8, for the AA, AB and AB' arrangements. Assume π-π stacking happens only between aromatic units. COF-9, obtained from COF-8, adsorbs UO2+ 2 ions and loses its fluorescence. In an experiment, a suspension of 5.000 mg COF-9 in 19.90 mg dm−3 UO2+ 2 solution (200.0 mL) was filtered after equilibrium was reached. The final concentration of UO2+ 2 was 9.225 mg dm−3.

\paragraph{Current subquestion.}
Calculate the equilibrium absorption capacity at a given temperature (𝑞𝑒) in mg g−1 of COF-9. Indicate the number of UO2+ 2 ions absorbed per pore. Assume that all uranium exists as UO2+ 2 ions.

\paragraph{Figure/image paths.}
\begin{itemize}
\item icho\_2026\_source/image/T3\_page-7.png
\item icho\_2026\_source/image/T3\_page-6.png
\end{itemize}

\paragraph{Formalization target.}
create a compiling Lean file with sorry bodies at `IChO2026Problems/problem\_icho\_2026\_t3\_a7.lean`.
The Lean declarations must preserve every mathematical object, quantifier, hypothesis, side condition, approximation/error guarantee, and requested conclusion from the source statement.
Use Mathlib/Physlib/CRNT names found through LeanExplore where available. Prefer the configured benchmark Base library; introduce a local definition only when the source genuinely requires an object that the environment does not expose.
This is an answer-blind solve-phase task. Derive a candidate result only from the problem-side evidence above; do not assume a target value, select a tolerance around a desired value, or encode a desired result into a candidate domain.
For numerical questions, define the raw end-to-end quantity and apply an explicit source-derived rounding rule. For identification questions, characterize or prove uniqueness before recording the derived candidate.

\begin{theorem}[Icho Chemistry formalization target]
\label{thm:physics:icho_2026_t3_a7:target}
The assigned autoformalize agent should translate this IChO chemistry problem into Lean declarations in the covered file, with theorem and lemma proof bodies written as `by sorry`.
\end{theorem}
\begin{proof}
This is an autoformalization task, not a proof task. Produce faithful statements that can later be proved without weakening the source contract.
\end{proof}
% --- Archon physics formalization source end ---
